\section{The quadratic estimate and its complete kernel}
\label{sec:quadratic}

All contractions in this section, unless a different metric is indicated, use
$g_0$.  We write $\dd V_0$ and $\dd A_0$ for its volume form and its induced
area form.  A function appearing in an integral over $S$ is understood to be
restricted to $S$.  Introduce
\[
 Q(v)=\Hess_{g_0}v-vB.
\]
The boundary normal $n=\partial_s$ points into $M$; the outward normal of
the domain $M$ along its inner boundary is therefore $-n$.

\Needspace{9\baselineskip}
\begin{lemma}[The weighted TT square identity]
\label{lem:tt-square}
Let $k$ be a smooth TT tensor of finite angular type, with decay
$k=O(r^{-\gamma})$ for some $\gamma>3/2$.  Let $v$ be the decaying solution of
\begin{equation}
 Lv=0,
 \qquad v|_S=D^{-1}(k_{ss}|_S).
 \label{eq:tt-dirichlet}
\end{equation}
Then $Q(v)$ is TT, $Q(v)_{ss}|_S=k_{ss}|_S$, and
\begin{equation}
 \int_M N|k|^2\,\dd V_0
 -\kappa\int_S vDv\,\dd A_0
 =\int_M N|k-Q(v)|^2\,\dd V_0.
 \label{eq:tt-square}
\end{equation}
All integrals in this identity converge.  In particular, no sign assumption
on $k_{ss}|_S$ is needed.
\end{lemma}

\begin{proof}
The operator $D=-\Delta_S+2\kappa/a$ is strictly positive on all spherical
harmonics.  Lemma~\ref{lem:inverse}, applied to the finitely many modes of
the boundary datum, gives the unique solution of
\eqref{eq:tt-dirichlet}; it has $v=O(r^{-3})$ with the differentiated
asymptotics needed below.

The background identities $\tr B=3$, $\Div B=0$, and
$B=\Ric_{g_0}+3g_0$ imply, by commuting the derivatives of a function,
\[
 \tr Q(v)=\Delta v-3v=Lv,
 \qquad
 (\Div Q(v))_i
 =\nabla_i\Delta v+(\Ric_{ij}-B_{ij})\nabla^jv
 =\nabla_iLv.
\]
Thus $Q(v)$ is TT.  Since $S$ is totally geodesic, the tangential trace of
$\Hess v$ on $S$ is $\Delta_Sv$.  Moreover $B_t(0)=\kappa/a$.
Taking the tangential trace of $Q(v)$ and using its vanishing full trace
therefore yields
\begin{equation}
 Q(v)_{ss}|_S=-\Delta_Sv+2(\kappa/a)v=Dv.
 \label{eq:q-normal-boundary}
\end{equation}

Set $\omega=N\,\dd v-v\,\dd N$.  With the averaged convention for
the symmetric gradient, $\operatorname{sym}\nabla\omega
=(\nabla_i\omega_j+\nabla_j\omega_i)/2$, the mixed products cancel and
the static identity $\Hess N=NB$ gives
\[
 \operatorname{sym}\nabla\omega
 =N\Hess v-v\Hess N=NQ(v).
\]
Symmetry and vanishing divergence of $k$ consequently give, on
$M_R=\{0\leq s\leq s(R)\}$,
\[
 \int_{M_R}N\langle k,Q(v)\rangle\,\dd V_0
 =\int_{\partial M_R}k(\nu,\omega^\sharp)\,\dd A_0.
\]
On $S$, $N=0$ and $\dd N=\kappa\,\dd s$, so that
$\omega^\sharp=-\kappa vn$.  Since $\nu=-n$ there, the inner contribution
is $+\kappa\int_Svk_{ss}\,\dd A_0$.  At infinity,
$|\omega|=O(r^{-2})$, and hence the outer contribution is
$O(r^{-\gamma})$ and tends to zero.  We obtain
\begin{equation}
 \int_M N\langle k,Q(v)\rangle\,\dd V_0
 =\kappa\int_Svk_{ss}\,\dd A_0
 =\kappa\int_SvDv\,\dd A_0.
 \label{eq:tt-pairing}
\end{equation}
The same calculation with $k$ replaced by $Q(v)$ gives
$\int_MN|Q(v)|^2\,\dd V_0=\kappa\int_SvDv\,\dd A_0$.
Expanding the square proves \eqref{eq:tt-square}.
Finally, $N\dd V_0$ has radial size $O(r^3)\,\dd s\,\dd A_\sigma$,
so the integral of $N|k|^2$ converges because $2\gamma>3$.
The stronger decay of $Q(v)$ gives all remaining convergence assertions.
\end{proof}

\subsection{The boundary area correction}

For the moment let $q$ have finite angular type, and use the conformal
metric jet $g=\phi^4g_0$ through order two.  Write $u=u_1[q]$, so that
$Lu=0$ and $u_s|_S=q_{ss}/4$.  Apply
Proposition~\ref{prop:cmc-jets}, and let $A_{\mathrm{leaf}}$ denote the
area of its first, minimal Taylor leaf.  This notation concerns that
formal leaf alone; it is not an enclosing-area infimum.  Let
\[
 U=[\II^{\circ}]_1,
 \qquad F=[f-\bar f]_1,
\]
where $f$ is the lapse and $\bar f$ its area average on each leaf.
Thus $F$ has zero round mean on each background sphere.

Let $v$ solve \eqref{eq:tt-dirichlet} with $k=q$.  At the fixed boundary,
the conformal boundary equation gives
\[
 H_g(S)=\eps\phi^{-6}q_{ss}.
\]
If $h_1$ is the first normal height of the minimal Taylor leaf, its
mean-curvature equation is
\[
 Dh_1+q_{ss}=0,
 \qquad h_1=-v|_S.
\]
The height difference from that leaf to $S$ is therefore
$\eps v|_S+O(\eps^2)$.  Area has zero first variation at a minimal
leaf, and its Hessian at the base sphere is $D$.  Taylor expansion about
the minimal leaf gives
\begin{equation}
 A(S)-A_{\mathrm{leaf}}
 =\frac{\eps^2}{2}\int_SvDv\,\dd A_0+O(\eps^3)
 \quad\text{as an identity of jets}.
 \label{eq:quadratic-area-gap}
\end{equation}
Indeed, a first-order perturbation of the area Hessian would multiply
two first-order displacements and first enter order three.  The first
variation at the minimal Taylor leaf vanishes to the order needed.
This reasoning also applies when its leading displacement points across
the boundary: it uses only the smooth boundary jets and their Taylor
extension, as in the construction of the formal leaves.

\begin{proposition}[The quadratic deficit]
\label{prop:quadratic}
On the real linear space of all smooth TT seeds with the prescribed
decay, the first deficit coefficient is zero and the second coefficient
$c_2[q]$ is a continuous, rotation-invariant, nonnegative quadratic form.
For a finite-angular-type seed it is given by
\begin{equation}
 \begin{gathered}
 c_2[q]=\frac1{16\pi}\int_M N\left(
 |q-Q(v)|^2+|U|^2+2FJ_s^0F\right)\,\dd V_0,\\
 J_s^0=-\Delta_{r^2\sigma}-\frac2{r^2}+\frac{6m_0}{r^3}.
 \end{gathered}
 \label{eq:quadratic-positive-form}
\end{equation}
The last term is integrated on the spheres before radial integration.
For every actual branch in Theorem~\ref{thm:main},
\begin{equation}
 \mathfrak D_q(\eps)=c_2[q]\eps^2+o(\eps^2).
 \label{eq:actual-quadratic-deficit}
\end{equation}
In particular, $c_2[q]>0$ implies a strictly positive deficit for all
sufficiently small positive $\eps$.
\end{proposition}

\begin{proof}
For a finite-angular-type seed the scalar constraint gives
$\scal_g+6=\eps^2|q|^2+O(\eps^3)$.
Proposition~\ref{prop:leading-energy}, with $l=1$, shows that the
constant and first coefficients of the mass norm minus
$b_*(A_{\mathrm{leaf}})$ vanish, and that its second coefficient is
\[
 \frac1{16\pi}\int_M N\bigl(|q|^2+|U|^2+2FJ_s^0F\bigr)\,\dd V_0.
\]
The areas in \eqref{eq:quadratic-area-gap} agree through order one,
so the first coefficient of the original deficit is also zero.
Since $b_*'(4\pi a^2)=\kappa/(8\pi)$, the contribution to its second
coefficient from that area difference is
$-\kappa(16\pi)^{-1}\int_SvDv\,\dd A_0$.
Lemma~\ref{lem:tt-square} now proves
\eqref{eq:quadratic-positive-form} with precisely the displayed factor
and sign.

On a spherical harmonic with $-\Delta_\sigma$ eigenvalue $\lambda\geq2$,
the eigenvalue of $J_s^0$ is
\[
 \frac{\lambda-2}{r^2}+\frac{6m_0}{r^3}>0.
\]
Thus its quadratic form is strictly positive on mean-zero functions at
every finite $s$, and $N>0$ for $s>0$.
This proves nonnegativity for finite-angular-type seeds without any
boundary sign restriction.

Proposition~\ref{prop:expansions} defines \(c_2\) on the full real
space \(\mathscr T_{\tau,\alpha}\) of smooth TT seeds as a continuous,
rotation-invariant quadratic form, independently of branch existence
and boundary sign, and proves \(c_1=0\). For an arbitrary seed \(q\),
take the finite-angular-type TT approximants \(q_\nu\) from
Lemma~\ref{lem:rotation}. The formula proved above gives
\(c_2[q_\nu]\geq0\), and continuity in \(\|\cdot\|_\beta\) gives
\(c_2[q]\geq0\). Proposition~\ref{prop:expansions} also supplies
\eqref{eq:actual-quadratic-deficit} for every prescribed actual branch.
\end{proof}

\subsection{All null directions}

Define the fixed vector space
\begin{equation}
 \mathcal K=\left\{Q(v):Lv=0,\ v\text{ decays},\quad
 v|_S\in\operatorname{span}_{\mathbb R}\{1,x_1,x_2,x_3\}\right\}.
 \label{eq:kernel-space}
\end{equation}
The decaying solutions in this definition have finite angular type and
$O(r^{-3})$ decay with all the differentiated bounds obtained from
their homogeneous radial equations.

\begin{corollary}[The complete quadratic kernel]
\label{cor:kernel}
The kernel of $c_2$ is exactly $\mathcal K$, a four-dimensional real
vector space.  In particular, the nonzero degree-one directions are
genuine quadratic null directions.
\end{corollary}

\begin{proof}
First suppose $q$ has finite angular type and $c_2[q]=0$.
Equation~\eqref{eq:quadratic-positive-form} and the strict positivity
just proved imply
\begin{equation}
 q=Q(v),\qquad U=0,\qquad F=0
 \quad\text{on }s>0.
 \label{eq:quadratic-equality-conditions}
\end{equation}
Smoothness extends these equalities to the boundary.  Coordinate
spheres are umbilic under conformal changes of $g_0$.  At the base
minimal sphere, the tracefree second-form variation due to a normal
height $h_1$ is $-(\Hess_Sh_1)^{\circ}$: the radial curvature term is
pure trace.  Since $h_1=-v|_S$, the vanishing of $U|_S$ implies
$(\Hess_\sigma(v|_S))^{\circ}=0$.
For completeness, the divergence of this equation on the unit sphere
gives $\dd(\Delta_\sigma v+2v)=0$.  Subtracting the mean reduces it to
$\Hess_\sigma(v-\bar v)=-(v-\bar v)\sigma$.
This equation determines its solution along every geodesic from its
value and differential at one point, and its solutions are precisely
the linear coordinate functions.  Thus $v|_S$ is a constant plus a
linear harmonic.

Uniqueness of the decaying Dirichlet problem for $L$ then excludes
every other angular mode throughout $M$.  Each degree-one radial
coefficient is a multiple of the same decaying solution of
\begin{equation}
 V''+2\frac NrV'-\left(3+\frac2{r^2}\right)V=0.
 \label{eq:dipole-radial-ode}
\end{equation}
Hence $q\in\mathcal K$.  Conversely, every element of $\mathcal K$
is null, as we now check; the conditions $U=F=0$ do not impose a further
restriction on its degree-one part.

Take $q=Q(v)\in\mathcal K$.  The first conformal coefficient $u$
also has only constant and linear modes, by its Neumann equation and
decaying uniqueness.  Let $h(s,x)$ be the first height coefficient
of the CMC spheres, and put $p=N/r$.  Conformal variation of mean
curvature and first variation of the lapse give
\begin{equation}
 [H]_1=J_s^0h+4(u_s-pu),
 \qquad [f]_1=h_s+2u.
 \label{eq:first-cmc-conformal-variations}
\end{equation}
Here the second identity follows by taking the normal component of
the velocity of the graph $s+\eps h(s,x)$ in
$g_0+4\eps u g_0$.
On the degree-one part, the CMC equation therefore fixes
\begin{equation}
 h=-\frac{2r^3}{3m_0}(u_s-pu).
 \label{eq:dipole-cmc-height}
\end{equation}
All formulas in the next calculation are componentwise on that part.
The background equations and \eqref{eq:dipole-radial-ode} give
\[
 p'=\frac{3m_0}{r^3}-\frac1{r^2},\qquad
 p^2=1+\frac1{r^2}-\frac{2m_0}{r^3},\qquad
 u''=-2pu'+\left(3+\frac2{r^2}\right)u.
\]
Consequently,
\begin{align}
 \frac1{r^3}\partial_s\bigl[r^3(u'-pu)\bigr]
 &=3p(u'-pu)+u''-p'u-pu'\notag\\
 &=\left(3+\frac2{r^2}-p'-3p^2\right)u
 =\frac{3m_0}{r^3}u.
 \label{eq:dipole-cancellation}
\end{align}
It follows that $h_s=-2u$ on the degree-one part, so its lapse
fluctuation vanishes.  The freely chosen radial mean of $h$ changes
$[f]_1$ only by a radial function, which disappears upon subtracting
the mean.  Thus $F=0$ independently of the radial labeling.

At every background sphere, the tracefree second-form variation of
these graphs is $-(\Hess_{r^2\sigma}h)^{\circ}$; the conformal
variation and radial curvature terms are pure trace.
The constant and linear modes of $h$ have vanishing tracefree
spherical Hessian, so $U=0$ as well.
At $S$ there is also exact agreement with the prescribed first minimal
height: on degree one,
\[
 D_1=\frac{6m_0}{a^3},\qquad
 h(0)=-\frac{4u_s(0)}{D_1}=-v(0).
\]
We already have $q=Q(v)$, and hence
\eqref{eq:quadratic-positive-form} proves $c_2[q]=0$.
Existence and uniqueness for each of the four boundary modes show
that $\mathcal K$ has dimension at most four; it has dimension
exactly four because \eqref{eq:q-normal-boundary} and invertibility
of $D$ make the map from these boundary data to $Q(v)$ injective.

It remains to rule out additional null directions of arbitrary angular
type.  Let $\mathcal Z=\{q:c_2[q]=0\}$ in the full smooth TT space.
A nonnegative quadratic form has a linear kernel: if $c_2[q]=0$,
nonnegativity of $c_2[q+th]$ for all real $t$ forces its polarized
bilinear form to vanish on $(q,h)$ for every $h$.
By Proposition~\ref{prop:quadratic}, $\mathcal Z$ is also closed in
$\|\cdot\|_\beta$ and invariant under rotations.
For $q\in\mathcal Z$, let $q_j$ be its polynomial rotational averages
from Lemma~\ref{lem:rotation}.  Each such average is a norm limit of
finite linear combinations of rotations of $q$; therefore
$q_j\in\mathcal Z$.  Each $q_j$ has finite angular type, so the
preceding argument puts it in $\mathcal K$.
Since $q_j\to q$ in norm and the fixed finite-dimensional space
$\mathcal K$ is closed in that norm, $q\in\mathcal K$.
This proves the full assertion without applying a pointwise geometric
limit to the CMC fields $U$ and $F$.
\end{proof}

\begin{remark}[The boundary sign and the null space]
\label{rem:kernel-sign}
If $v|_S=c+d\cdot x$, then
\[
 Q(v)_{ss}|_S=D_0c+D_1d\cdot x,
 \qquad D_0=\frac{2\kappa}{a},\quad
 D_1=\frac2{a^2}+\frac{2\kappa}{a}=\frac{6m_0}{a^3}.
\]
Thus $Q(v)_{ss}|_S\geq0$ is equivalent to $D_0c\geq D_1|d|$.
A nonzero pure dipole fails that condition, but a sufficiently large
positive radial component makes a mixed radial--dipole seed satisfy
it.  This is a compatibility statement about seeds, not an existence
claim for branches satisfying outermostness and outer
area-minimization.  In particular, the boundary sign cannot be used
to discard the degree-one null directions from the proof.
\end{remark}
